\tikzset{blk/.style={draw,minimum height=7mm,minimum width=9mm,fill=white,font=\small},
 sumj/.style={draw,circle,inner sep=1.2pt,font=\scriptsize},
 tap/.style={circle,fill,inner sep=1.3pt},
 ar/.style={-Stealth}}

\section{Chapter 1: Linear systems and their description \texorpdfstring{\fr{6--30}}{}}
\label{ch1}

\subsection*{Big picture}
A \textbf{linear system} is a machine with input $u(t)$, output $y(t)$ and internal \textbf{state} $x(t)$.
The state is ``everything you need to know now, so that together with future inputs you can predict the future''.
In this chapter we learn to \emph{draw} systems (block diagrams) and to \emph{rewrite} any constant-coefficient ODE as a first order
matrix system $\dot X=AX+Bu$. Nothing is solved yet; that is Chapter 2.

\subsection{Linear operators and the five building blocks \texorpdfstring{\fr{7--9}}{}}\label{ss:1-1}
\begin{plainbox}
A linear system is a \textbf{machine}: a signal $u(t)$ goes in, a signal $y(t)$ comes out. Inside, the signal passes through
simple ``boxes''. Each box does one job: \emph{scale it} (amplifier), \emph{add up the history} (integrator), \emph{measure how fast it changes}
(differentiator), \emph{combine signals} (summing junction) or \emph{copy it} (tap).
\textbf{Linear} means ``double the input, double the output; add two inputs, add the outputs''.
\par\smallskip\textbf{In the exam this looks like:} ``Show that $L$ is linear'' or ``which of these maps are linear?''.
\end{plainbox}
\begin{defbox}
$L:V\to W$ is \textbf{linear} if $L(\lambda f+\mu g)=\lambda Lf+\mu Lg$ for all $f,g$ and all $\lambda,\mu\in\CC$.
The basic linear operators (on continuously differentiable $f:[0,\infty)\to\CC$) are
\begin{enumerate}[label=(\roman*)]
\item \textbf{differentiator} $f\mapsto df/dt$, drawn $[d/dt]$;
\item \textbf{integrator} $f\mapsto\int_0^x f(t)\,dt$, drawn $[\int]$ (always integrate from $0$);
\item \textbf{amplifier} $f\mapsto af$ for a constant $a\in\CC$, drawn $[a]$ (called an amplifier even if $|a|<1$);
\item \textbf{multiplication} by a fixed function $h$: $f\mapsto hf$, drawn $[h]$;
\item \textbf{evaluation} $f\mapsto f(t_0)$, drawn $[\delta_{t_0}]$.
\end{enumerate}
Two connectors: a \textbf{summing junction} $\oplus$ adds the incoming signals (with the signs written on the arrows);
a \textbf{tap} $\bullet$ copies a signal onto several wires \emph{without} dividing it (like a point in an electrical circuit,
all wires at a tap carry the same voltage).
\end{defbox}

\begin{warnbox}
A tap does \textbf{not} split the signal into halves. If $u$ goes into a tap with three wires out, each wire carries $u$.
\end{warnbox}

\begin{exambox}
``Show that $L$ is linear'' questions: write out $L(\lambda f+\mu g)$ and use linearity of differentiation/integration.
E.g.\ $Lf=t\,f''+3\int_0^t f$ is linear since $L(\lambda f+\mu g)=t(\lambda f''+\mu g'')+3\int_0^t(\lambda f+\mu g)=\lambda Lf+\mu Lg$.
A map like $f\mapsto f^2$ or $f\mapsto f+1$ is \textbf{not} linear (check $L(2f)\neq 2Lf$, or $L(0)\neq0$).
\end{exambox}

\subsection{Block diagrams: rules, open loop and feedback \texorpdfstring{\fr{11--17}}{}}\label{ss:1-2}
\begin{plainbox}
\textbf{Feedback} means part of the output is sent back to the input. Think of a thermostat: the room temperature (output)
is measured and used to switch the heater (input). \textbf{Negative} feedback calms the system down (going faster creates drag);
\textbf{positive} feedback pushes it further the same way (more birds $\Rightarrow$ more births). Open loop = no signal ever comes back.
\par\smallskip\textbf{In the exam this looks like:} ``Find the output of this feedback diagram'' (write the equation at the summing junction and solve for $y$).
\end{plainbox}
\begin{defbox}
A block diagram is a directed graph from input $u$ to output $y$ built from $u,y,\delta_0,\oplus,\bullet,d/dt,\int,a,h$.
\begin{itemize}
\item $u$ and $y$ have degree one; every other vertex has degree two or three.
\item No loops at a vertex and no multiple edges; every vertex lies on a directed path from $u$ to $y$.
\item If the diagram contains a \textbf{circuit} there is \textbf{feedback} (closed loop); otherwise it is \textbf{open loop} (straight through).
\end{itemize}
\end{defbox}

\begin{workbox}[title={Worked example: Black's amplifier \fr{15} (positive feedback)}]
\begin{center}
\begin{tikzpicture}
\node (u) at (0,0) {$u$};
\node[sumj] (s) at (1.6,0) {$+$};
\node[tap] (t) at (4,0) {};
\node (y) at (5.3,0) {$y$};
\node[blk] (th) at (2.8,1.1) {$\theta$};
\draw[ar] (u)--(s); \draw[ar] (s)--node[above]{$u+\theta y$}(t); \draw[ar] (t)--(y);
\draw[ar] (t)|-(th); \draw[ar] (th)-|node[left,pos=.8]{$\theta y$}(s);
\end{tikzpicture}
\end{center}
\textbf{Step 1.} Write the equation at the summing junction: the signal leaving it is $u+\theta y$, and that is $y$.\\
\textbf{Step 2.} So $y=u+\theta y$, hence $(1-\theta)y=u$ and $\boxed{y=\dfrac{u}{1-\theta}}$ ($0<\theta<1$).\\
Since $1/(1-\theta)>1$ the input is amplified; but so is any noise, so it is not practical on its own.
Note: summing junction next to the input, tap next to the output.
\end{workbox}

\begin{workbox}[title={Worked example: negative feedback from air resistance \fr{16--17}}]
Cyclist: $m\,\dfrac{dv}{dt}=-kv+u$. Integrate from $0$ (with $v(0)=0$) and divide by $m$:
\[ v=-\frac km\int v+\frac1m\int u. \]
The \textbf{minus} sign of $-k/m$ on the fed-back signal is what ``negative feedback'' means: going faster produces a force that slows you down.
Positive feedback example: bird population $\dot x=kx+u$ with $k>0$ (more birds $\Rightarrow$ more births); turbocharger.
\end{workbox}

\subsection{How to turn an equation into a block diagram}\label{sec:bd}\label{ss:1-3}
\begin{plainbox}
Differentiators amplify noise, so we want diagrams built from \textbf{integrators} only.
The trick: \textbf{integrate the whole equation} until no derivatives are left. Then every term is ``an integral of $y$'' or
``an integral of $u$''. Terms with $y$ become a loop back from the output; terms with $u$ go forward from the input.
\par\smallskip\textbf{In the exam this looks like:} ``Express this ODE / integral equation as a block diagram with no differentiators'' (AE1, A1.1--A1.2). Usually 5 marks.
\end{plainbox}
\begin{methodbox}
\begin{enumerate}
\item If differentiators \emph{are allowed}: tap $u$ into one branch per term, put the right operator ($d/dt$, $\int$, amplifier) on each branch, and add the branches at a summing junction. This is open loop if $y$ does not appear on the right.
\item If \textbf{no differentiators} are allowed: \textbf{integrate both sides} enough times that no derivative is left
(as many times as the highest derivative of $y$ or $u$), assuming zero initial values.
\item Solve for $y$ (divide by the coefficient of $y$). You now have $y=(\text{multiple integrals of }y)+(\text{multiple integrals of }u)$.
\item Draw: terms involving $y$ are a \textbf{feedback} loop from a tap at the output back to a summing junction; terms involving $u$ feed forward.
\item Check: read the diagram backwards and recover the equation.
\end{enumerate}
\end{methodbox}

\begin{workbox}[title={Worked example (AE1, A1.1(i)): $y=4\dfrac{d^2u}{dt^2}+3\displaystyle\int u+6u$}]
Differentiators allowed and $y$ does not appear on the right, so this is \textbf{open loop}: tap $u$ three ways.
\begin{center}
\begin{tikzpicture}
\node (u) at (0,0) {$u$}; \node[tap] (t) at (1,0) {};
\node[blk] (d1) at (2.4,1.2) {$d/dt$}; \node[blk] (d2) at (3.8,1.2) {$d/dt$}; \node[blk] (a4) at (5.2,1.2) {$4$};
\node[blk] (i) at (2.4,0) {$\int$}; \node[blk] (a3) at (5.2,0) {$3$};
\node[blk] (a6) at (5.2,-1.2) {$6$};
\node[sumj] (s) at (7,0) {$+$}; \node (y) at (8.3,0) {$y$};
\draw[ar] (u)--(t); \draw[ar] (t)|-(d1); \draw[ar] (d1)--(d2); \draw[ar] (d2)--(a4);
\draw[ar] (t)--(i); \draw[ar] (i)--(a3); \draw[ar] (t)|-(a6);
\draw[ar] (a4)-|node[above,pos=.3]{$4u''$}(s); \draw[ar] (a3)--node[above]{$3\!\int\! u$}(s); \draw[ar] (a6)-|node[below,pos=.3]{$6u$}(s); \draw[ar] (s)--(y);
\end{tikzpicture}
\end{center}
\end{workbox}

\begin{workbox}[title={Worked example (AE1, A1.1(ii)): $y+4\dfrac{d^2y}{dt^2}=2\dfrac{du}{dt}+7\displaystyle\int u$, no differentiators}]
\textbf{Step 1: integrate twice} (highest derivative is $y''$), zero initial conditions, writing $\int^{(2)}$ for a double integral:
\[ \textstyle\int^{(2)}y+4y=2\int u+7\int^{(3)}u. \]
\textbf{Step 2: solve for $y$:}
\[ \boxed{\,y=-\tfrac14\textstyle\int^{(2)}y+\tfrac12\int u+\tfrac74\int^{(3)}u\,} \]
\textbf{Step 3: nest the integrals} (a neat way to use only three integrators):
\[ y=\int\Big(\tfrac12u+\int\big(-\tfrac14 y+\int\tfrac74u\big)\Big). \]
(Check by expanding: $\int\frac12 u-\frac14\int^{(2)}y+\frac74\int^{(3)}u$. \checkmark)
\begin{center}
\begin{tikzpicture}
\node (u) at (-0.6,0) {$u$}; \node[tap] (t) at (0.3,0) {};
\node[blk] (a) at (1.3,0) {$7/4$}; \node[blk] (i1) at (2.7,0) {$\int$};
\node[sumj] (s1) at (3.9,0) {$+$}; \node[blk] (i2) at (5.0,0) {$\int$};
\node[sumj] (s2) at (6.2,0) {$+$}; \node[blk] (i3) at (7.3,0) {$\int$};
\node[tap] (ty) at (8.5,0) {}; \node (y) at (9.4,0) {$y$};
\node[blk] (h) at (6.2,1.3) {$1/2$};
\node[blk] (f) at (6.2,-1.3) {$-1/4$};
\draw[ar] (u)--(t); \draw[ar] (t)--(a); \draw[ar] (a)--(i1); \draw[ar] (i1)--(s1); \draw[ar] (s1)--(i2); \draw[ar] (i2)--(s2);
\draw[ar] (s2)--(i3); \draw[ar] (i3)--(ty); \draw[ar] (ty)--(y);
\draw[ar] (t)|-(h); \draw[ar] (h)--(s2);
\draw[ar] (ty)|-(f); \draw[ar] (f)-|(s1);
\end{tikzpicture}
\end{center}
The loop $y\to[-1/4]\to\oplus\to\int\to\oplus\to\int\to y$ is the feedback; there are no differentiators.
\end{workbox}

\begin{workbox}[title={Worked example (AE1, A1.2): coupled equations $\dot x=ax+bu,\ \dot y=cx+du$}]
\textbf{Step 1.} Integrate each (zero initial values): $x=a\int x+b\int u$ and $y=\int(cx+du)$.\\
\textbf{Step 2.} Draw $x$ as a feedback loop around an integrator, then feed $x$ and $u$ forward.
\begin{center}
\begin{tikzpicture}
\node (u) at (-0.5,0) {$u$}; \node[tap] (t) at (0.4,0) {};
\node[blk] (b) at (1.4,0) {$b$}; \node[sumj] (s1) at (2.6,0) {$+$}; \node[blk] (i1) at (3.7,0) {$\int$};
\node[tap] (tx) at (4.8,0) {}; \node[blk] (a) at (3.7,1.2) {$a$};
\node[blk] (c) at (5.8,0) {$c$}; \node[sumj] (s2) at (7,0) {$+$}; \node[blk] (i2) at (8.1,0) {$\int$}; \node (y) at (9.2,0) {$y$};
\node[blk] (d) at (5.8,-1.2) {$d$};
\draw[ar] (u)--(t); \draw[ar] (t)--(b); \draw[ar] (b)--(s1); \draw[ar] (s1)--(i1); \draw[ar] (i1)--node[above]{$x$}(tx);
\draw[ar] (tx)|-(a); \draw[ar] (a)-|(s1); \draw[ar] (tx)--(c); \draw[ar] (c)--(s2); \draw[ar] (s2)--(i2); \draw[ar] (i2)--(y);
\draw[ar] (t)|-(d); \draw[ar] (d)-|(s2);
\end{tikzpicture}
\end{center}
\textbf{(ii) Two inputs.} With $\dot x=ax+b_1u_1+b_2u_2$, $\dot y=cx+d_1u_1+d_2u_2$: the same diagram, but the first summing junction
receives $b_1u_1$ and $b_2u_2$ (two amplifiers) and the second receives $cx$, $d_1u_1$, $d_2u_2$. Equivalently use the vector input $U=[u_1;u_2]$
and matrix amplifiers $[b_1\ b_2]$ and $[d_1\ d_2]$ \fr{18}.
\end{workbox}

\begin{thmbox}[title={Proposition 0.1 \fr{20}: ODEs are feedback systems}]
$a_ny^{(n)}+\dots+a_0y=b_mu^{(m)}+\dots+b_0u$ (constant coefficients, $a_n\neq0$) can be realised with taps, amplifiers, summing junctions,
integrators and differentiators. \textbf{Idea of proof:} integrate $n$ times and divide by $a_n$:
\[ y=-\frac{a_{n-1}}{a_n}\int y-\dots-\frac{a_0}{a_n}\int^{(n)}y+\frac{b_m}{a_n}\int^{(n)}u^{(m)}+\dots+\frac{b_0}{a_n}\int^{(n)}u ,\]
build $x=\frac{b_m}{a_n}u^{(m)}+\dots+\frac{b_0}{a_n}u$ with differentiators (input operator, \fr{21}), and feed $y$ back through integrators (\fr{22}).
\end{thmbox}

\subsection{Reduction of order and the matrix form \texorpdfstring{\fr{23--29}}{}}\label{sec:reduction}\label{ss:1-4}
\begin{plainbox}
A big $n$-th order ODE is hard; $n$ small first-order ODEs are easy, and they fit into \textbf{one matrix equation}.
So we keep track of the list $y, y', \dots, y^{(n-1)}$ as one vector $X$ (the \emph{state}).
Each entry's derivative is simply the next entry, except the last, which the ODE tells us.
\par\smallskip\textbf{In the exam this looks like:} ``Introduce $v=dx/dt$ and write the system as $(A,B,C,D)$'' (AE1, A1.3).
\end{plainbox}
\begin{thmbox}[title={Lemma 0.2 (reduction of order) \fr{26--27}}]
$y^{(n)}+a_{n-1}(t)y^{(n-1)}+\dots+a_0(t)y=u(t)$ is equivalent to $\dfrac{dX}{dt}=AX+Bu$ with
\[
X=\begin{bmatrix}y\\y'\\\vdots\\y^{(n-1)}\end{bmatrix},\quad
A=\begin{bmatrix}0&1&0&\cdots&0\\0&0&1&\cdots&0\\\vdots&&&\ddots&\vdots\\0&0&\cdots&0&1\\-a_0&-a_1&\cdots&\cdots&-a_{n-1}\end{bmatrix},\quad
B=\begin{bmatrix}0\\0\\\vdots\\0\\1\end{bmatrix}.
\]
\textbf{Why:} the first $n-1$ rows just say ``the derivative of $y^{(j)}$ is $y^{(j+1)}$''; the last row is the ODE solved for $y^{(n)}$.
\end{thmbox}

\begin{methodbox}[title={Method: write an ODE as $(A,B,C,D)$}]
\begin{enumerate}
\item Make the ODE \textbf{monic}: divide by the leading coefficient.
\item Choose states $x_1=y,\ x_2=y',\dots,x_n=y^{(n-1)}$.
\item Write $\dot x_1=x_2,\ \dots,\ \dot x_{n-1}=x_n,\ \dot x_n=(\text{ODE solved for }y^{(n)})$.
\item Read off $A$ and $B$. The output is $y=x_1$, so $C=[1\ 0\ \cdots\ 0]$, $D=0$.
\end{enumerate}
\end{methodbox}

\begin{workbox}[title={Worked example (AE1, A1.3): simple harmonic oscillator $m\ddot x+kx=u$}]
\textbf{Step 1.} Monic: $\ddot x=-\frac km x+\frac1m u$.\quad
\textbf{Step 2.} States $x$ and $v=\dot x$.\quad
\textbf{Step 3.} $\dot x=v$, $\dot v=-\frac km x+\frac1m u$.\\
\textbf{Step 4.}
\[
\frac{d}{dt}\begin{bmatrix}x\\v\end{bmatrix}=\underbrace{\begin{bmatrix}0&1\\-k/m&0\end{bmatrix}}_{A}\begin{bmatrix}x\\v\end{bmatrix}+\underbrace{\begin{bmatrix}0\\1/m\end{bmatrix}}_{B}u,\qquad
x=\underbrace{\begin{bmatrix}1&0\end{bmatrix}}_{C}\begin{bmatrix}x\\v\end{bmatrix}+\underbrace{0}_{D}\,u.
\]
\end{workbox}

\begin{workbox}[title={Worked example: damped harmonic oscillator \fr{23--25}}]
$y''+\beta y'+\gamma y=u$ (car suspension: spring $-\gamma y$, shock absorber $-\beta y'$). With $v=y'$:
$A=\begin{bmatrix}0&1\\-\gamma&-\beta\end{bmatrix}$, $B=\begin{bmatrix}0\\1\end{bmatrix}$, $C=[1\ 0]$, $D=0$.
As a feedback system: $y=-\beta\int y-\gamma\int^{(2)}y+\int^{(2)}u$. This example comes back in Chapters 2, 4 and 5.
\end{workbox}

\subsection{The companion matrix \texorpdfstring{\fr{30}, \fr{34--35}}{}}\label{sec:companion}\label{ss:1-5}
\begin{plainbox}
The companion matrix is a \textbf{ready-made matrix for any polynomial}: put ones just above the diagonal and the
polynomial's coefficients (with minus signs, constant term first) along the bottom row. Its characteristic polynomial is exactly
the polynomial you started with.
\par\smallskip\textbf{In the exam this looks like:} ``Find a matrix whose characteristic polynomial is $p(s)$'' (AE2, A2.2), or the induction proof.
\end{plainbox}
\begin{thmbox}[title={Proposition 0.5: every monic polynomial is a characteristic polynomial}]
The companion matrix $A_c$ (last row $-a_0,-a_1,\dots,-a_{n-1}$, ones on the superdiagonal) has
\[ \det(\lambda I-A_c)=\lambda^n+a_{n-1}\lambda^{n-1}+\dots+a_1\lambda+a_0. \]
\end{thmbox}
\begin{warnbox}
Note the \textbf{order} and \textbf{signs}: the last row is $[-a_0\ \ -a_1\ \cdots\ -a_{n-1}]$, i.e.\ the coefficients from the
\emph{constant term upwards}, all with a \emph{minus} sign.
\end{warnbox}
\begin{proofbox}[title={Proof of Proposition 0.5 (induction on $n$; expand along the first column)}]
$P_1$ is trivial. Assume true for size $n-1$. Expanding $\det(\lambda I-A_c)$ along the first column, which has $\lambda$ in place $(1,1)$ and $a_0$ in place $(n,1)$:
\[ \det(\lambda I-A_c)=\lambda\det\begin{bmatrix}\lambda&-1&&\\&\ddots&\ddots&\\&&\lambda&-1\\a_1&\cdots&a_{n-2}&\lambda+a_{n-1}\end{bmatrix}+(-1)^{n+1}a_0\det\begin{bmatrix}-1&&&\\\lambda&-1&&\\&\ddots&\ddots&\\&&\lambda&-1\end{bmatrix}. \]
The first determinant is $\lambda^{n-1}+a_{n-1}\lambda^{n-2}+\dots+a_1$ by the induction hypothesis (it is the companion form for $a_1,\dots,a_{n-1}$).
The second matrix is lower triangular with $n-1$ diagonal entries $-1$, so its determinant is $(-1)^{n-1}$, and $(-1)^{n+1}(-1)^{n-1}=1$. Hence
$\det(\lambda I-A_c)=\lambda(\lambda^{n-1}+\dots+a_1)+a_0=\lambda^n+a_{n-1}\lambda^{n-1}+\dots+a_0$.
\end{proofbox}

\begin{workbox}[title={Worked example (AE2, A2.2): matrix with $\chi_A(s)=s^4+2s^3+s^2+4s+2$}]
Here $a_3=2,\ a_2=1,\ a_1=4,\ a_0=2$, so
\[ A=\begin{bmatrix}0&1&0&0\\0&0&1&0\\0&0&0&1\\-2&-4&-1&-2\end{bmatrix}. \]
Numerically (MATLAB \texttt{eig(A)} or \texttt{roots([1 2 1 4 2])}): $-2.1877,\ -0.5156,\ 0.3516\pm1.2843i$.\\
\textit{Bonus:} two eigenvalues have positive real part, so $\dot X=AX$ is \textbf{unstable}. (Coefficients all positive does \emph{not} imply stable, see \S\ref{sec:stabpoly}.)
\end{workbox}

\subsection{Exam-style questions for Chapter 1}\label{ss:1-6}
\begin{answerbox}
\textbf{Block diagram:} (1) integrate $k$ times, saying ``with zero initial conditions''; (2) solve for $y$; (3) draw, labelling every wire; (4) one line checking the diagram gives back the equation.\\
\textbf{Matrix form:} (1) make monic; (2) define the state vector in words (``let $x_1=y$, $x_2=y'$''); (3) write the first-order equations; (4) state $A,B,C,D$ clearly with their sizes.
\end{answerbox}
\begin{qbox}[title={Q1.1 (easy)}]
Express $\dfrac{d^2y}{dt^2}+3\dfrac{dy}{dt}+2y=5u$ (a) as a block diagram without differentiators and (b) as an $(A,B,C,D)$ system.
\end{qbox}
\begin{solbox}
(a) Integrate twice: $y+3\int y+2\int^{(2)}y=5\int^{(2)}u$, so $y=-3\int y-2\int^{(2)}y+5\int^{(2)}u$.
Nested: $y=\int\big(-3y+\int(5u-2y)\big)$. Diagram: $u\to[5]\to\oplus\to\int\to\oplus\to\int\to\bullet\to y$, with $y$ fed back through $[-2]$ into the first $\oplus$ and through $[-3]$ into the second.\\
(b) $X=[y;\,y']$, $A=\begin{bmatrix}0&1\\-2&-3\end{bmatrix}$, $B=\begin{bmatrix}0\\5\end{bmatrix}$, $C=[1\ 0]$, $D=0$.
(Equally valid: $B=[0;1]$, $C=[5\ 0]$. Realisations are not unique.)
\end{solbox}

\begin{qbox}[title={Q1.2 (medium)}]
Write $y'''-2y''+y'+7y=u$ in the form $\dot X=AX+Bu$, $y=CX$, and state $\det(sI-A)$ without computing a determinant.
\end{qbox}
\begin{solbox}
$X=[y;y';y'']$, $A=\begin{bmatrix}0&1&0\\0&0&1\\-7&-1&2\end{bmatrix}$, $B=[0;0;1]$, $C=[1\ 0\ 0]$.
$A$ is a companion matrix with $a_0=7,a_1=1,a_2=-2$, so by Proposition 0.5, $\det(sI-A)=s^3-2s^2+s+7$.
\end{solbox}

\begin{qbox}[title={Q1.3 (harder: reading a diagram)}]
In a diagram, $u$ enters a summing junction; the junction output goes through an integrator to a tap; the tap gives $y$, and also feeds $y$ back through
an amplifier $[-3]$ into the summing junction and forward through $[2]$ to a second junction that also receives $u$. The output of the second junction is $z$. Find the ODE for $y$ and express $z$ in terms of $u$ and $y$.
\end{qbox}
\begin{solbox}
The first junction outputs $u-3y$, and integrating gives $y$: $y=\int(u-3y)$, so $\dot y=-3y+u$. The second junction: $z=2y+u$.
This is the SISO system $A=[-3]$, $B=[1]$, $C=[2]$, $D=[1]$.
\end{solbox}
\newpage
